% ============================================================
% Academic Poster: Automated Screening for Cognitive Decline
% Author: Feixiang Xiao
% Course: Introduction to Deep Learning — Bonus Seminar
% University of Stuttgart, 2025
% Requires: beamerposter, tikz, booktabs, multicol, colortbl
% Compile: pdflatex poster_cognitive_decline.tex (run twice)
% ============================================================

\documentclass[final]{beamer}

\usepackage[size=a0,orientation=portrait,scale=1.35]{beamerposter}
\usepackage{graphicx}
\usepackage{booktabs}
\usepackage{array}
\usepackage{tikz}
\usetikzlibrary{arrows.meta,positioning,shapes.geometric,fit,backgrounds}
\usepackage[T1]{fontenc}
\usepackage[utf8]{inputenc}
\usepackage{lmodern}
\usepackage{microtype}
\usepackage{amsmath}
\usepackage{xcolor}
\usepackage{colortbl}
\usepackage{url}
\usepackage{enumitem}
\usepackage{multicol}
\usepackage{ragged2e}
\usepackage{calc}

% ---- Color Palette ----------------------------------------
\definecolor{Cprimary}{RGB}{0,62,116}       % Stuttgart dark blue
\definecolor{Csecondary}{RGB}{0,100,170}    % Medium blue
\definecolor{Clight}{RGB}{218,235,248}      % Light blue tint
\definecolor{Cgreen}{RGB}{0,110,84}         % Dark green
\definecolor{Corange}{RGB}{190,75,0}        % Warm orange (alert)
\definecolor{Cdark}{RGB}{20,20,30}          % Near-black body text
\definecolor{Crule}{RGB}{180,200,220}       % Table rule color

% ---- Beamer Theme Setup -----------------------------------
\usetheme{default}
\useinnertheme{rectangles}
\setbeamertemplate{navigation symbols}{}
\setbeamertemplate{footline}{}

\setbeamercolor{block title}{bg=Csecondary, fg=white}
\setbeamercolor{block body}{bg=Clight, fg=Cdark}
\setbeamercolor{block alerted title}{bg=Corange, fg=white}
\setbeamercolor{block alerted body}{bg=orange!10, fg=Cdark}
\setbeamercolor{block example title}{bg=Cgreen, fg=white}
\setbeamercolor{block example body}{bg=green!9, fg=Cdark}

\setbeamerfont{block title}{size=\large, series=\bfseries}
\setbeamerfont{block body}{size=\normalsize}

% ---- Headline (custom header) -----------------------------
\setbeamertemplate{headline}{%
  \begin{beamercolorbox}[wd=\paperwidth, colsep=0pt]{headline}%
    \begin{tikzpicture}
      \fill[Cprimary] (0,0) rectangle (\paperwidth, 13cm);
      \fill[Csecondary!60] (0,0) rectangle (\paperwidth, 0.5cm);
    \end{tikzpicture}%
    \vspace{-13.5cm}%
    \begin{center}
      {\fontsize{56}{66}\selectfont\bfseries\color{white}%
        Automated Screening for Cognitive Decline}\\[0.45em]
      {\fontsize{44}{52}\selectfont\bfseries\color{white!85}%
        From Picture Descriptions to Algorithms}\\[1.1em]
      {\fontsize{30}{38}\selectfont\color{white!80}%
        \textbf{Feixiang Xiao}\enspace\textbar\enspace
        Introduction to Deep Learning — Bonus Seminar\enspace\textbar\enspace
        University of Stuttgart, 2025}\\[1.2em]
    \end{center}
  \end{beamercolorbox}%
}

% ---- Layout dimensions ------------------------------------
\newlength{\postermargin}
\setlength{\postermargin}{1.8cm}
\newlength{\colsep}
\setlength{\colsep}{1.4cm}
\newlength{\leftcol}
\newlength{\rightcol}
\setlength{\leftcol}{0.475\linewidth}
\setlength{\rightcol}{0.475\linewidth}
\newlength{\blockskip}
\setlength{\blockskip}{1.0cm}

% ===========================================================
\begin{document}
\begin{frame}[t]

\vspace{1.8cm}

% -----------------------------------------------------------
% TWO-COLUMN BODY
% -----------------------------------------------------------
\begin{columns}[T, totalwidth=\linewidth]

% ===========================================================
% LEFT COLUMN
% ===========================================================
\begin{column}{\leftcol}

% -----------------------------------------------------------
\begin{block}{1\enspace Motivation}
\justifying
\begin{itemize}[leftmargin=1.4em, itemsep=0.35em, topsep=0pt]
  \item \textbf{57 million} people worldwide live with dementia;
        Alzheimer's disease (AD) accounts for 60--70\% of cases
  \item Early detection enables timely intervention and reduces caregiver burden
  \item Standard tools (MMSE, PET/MRI) are \textbf{expensive,
        clinician-dependent}, and inaccessible at population scale
  \item \textbf{Speech} is a non-invasive, low-cost, remotely deployable biomarker
        that reflects language and cognition simultaneously
\end{itemize}
\vspace{0.6em}
\begin{center}
\begin{tikzpicture}
  \node[
    draw=Cprimary, fill=Cprimary!12, rounded corners=7pt,
    inner sep=12pt, text width=0.86\textwidth, align=justify
  ] {
    \textbf{Research Question:} How well can automated speech-based
    systems screen for cognitive decline --- and what separates
    benchmark accuracy from real-world clinical deployment?
  };
\end{tikzpicture}
\end{center}
\end{block}

\vspace{\blockskip}

% -----------------------------------------------------------
\begin{block}{2\enspace The Cookie Theft Task}
\justifying
Introduced in the \textit{Boston Diagnostic Aphasia Examination}
(Goodglass \& Kaplan, 1983), the \textbf{Cookie Theft picture
description task} asks participants to describe a scene in which:
a child steals cookies from a jar while standing on a tipping stool;
a woman washes dishes, oblivious to the overflowing sink.

\vspace{0.5em}
\noindent The task elicits \textbf{spontaneous connected speech},
revealing characteristic patterns of cognitive decline:

\begin{itemize}[leftmargin=1.4em, itemsep=0.25em, topsep=0.3em]
  \item Word-finding difficulties and circumlocutions
  \item Reduced syntactic complexity and sentence length
  \item Missing information units (IUs) and topic repetition
  \item Increased disfluency, hesitations, and pause frequency
\end{itemize}

\vspace{0.6em}
\begin{center}
\renewcommand{\arraystretch}{1.35}
\begin{tabular}{@{}p{0.38\textwidth} p{0.56\textwidth}@{}}
  \toprule
  \rowcolor{Cprimary!15}
  \textbf{Resource} & \textbf{Description} \\
  \midrule
  DementiaBank / Pitt Corpus
    & 208 AD + 104 HC; longitudinal recordings
      (Becker et al., 1994) \\
  ADReSS Challenge
    & 78 AD + 78 HC; age/gender-balanced,
      acoustically pre-processed benchmark
      (Luz et al., 2020) \\
  ADReSSo Challenge
    & Speech-only variant; adds longitudinal
      decline prediction task
      (Luz et al., 2021) \\
  ADReSS-M (ICASSP 2023)
    & Cross-lingual: trained on English,
      tested on Greek
      (Luz et al., 2024) \\
  \bottomrule
\end{tabular}
\end{center}
\end{block}

\vspace{\blockskip}

% -----------------------------------------------------------
\begin{block}{3\enspace Processing Pipeline}
\vspace{0.4em}
\begin{center}
\begin{tikzpicture}[
  node distance = 0pt,
  box/.style = {
    draw = Csecondary, fill = Clight,
    rounded corners = 6pt,
    minimum width  = 5.2cm,
    minimum height = 1.6cm,
    text width     = 4.8cm,
    align = center,
    font = \small\bfseries\color{Cdark}
  },
  mainbox/.style = {
    box,
    fill  = Cprimary!18,
    draw  = Cprimary
  },
  outbox/.style = {
    box,
    fill  = Cgreen!20,
    draw  = Cgreen
  },
  featurebox/.style = {
    box,
    minimum width  = 4.6cm,
    text width     = 4.2cm
  },
  arr/.style = {
    -{Latex[length=9pt, width=6pt]},
    line width = 2.2pt,
    color = Csecondary
  },
  darr/.style = {
    arr,
    color = Cprimary
  },
  lbl/.style = {
    font = \scriptsize\itshape,
    text = Cdark,
    align = center
  }
]

% Nodes
\node[mainbox] (audio) {Speech\\Recording};

\node[mainbox, right=3.2cm of audio] (asr)
  {ASR /\\Feature Extraction};

\node[featurebox, above right=0.9cm and 2.8cm of asr] (acou)
  {Acoustic\\Features};

\node[featurebox, below right=0.9cm and 2.8cm of asr] (ling)
  {Linguistic\\Features};

\node[mainbox, right=2.8cm of asr, yshift=0cm] (clf)
  {ML Classifier};

\node[outbox, right=2.8cm of clf] (out)
  {Screening\\Decision};

% Arrows
\draw[darr] (audio) -- (asr);
\draw[darr] (asr.east) -- ++(1.2,0) |- (acou.west);
\draw[darr] (asr.east) -- ++(1.2,0) |- (ling.west);
\draw[darr] (acou.east) -| (clf.north);
\draw[darr] (ling.east) -| (clf.south);
\draw[darr] (clf) -- (out);

% Labels under feature nodes
\node[lbl, below=0.18cm of acou]
  {eGeMAPS, prosody,\\pauses, disfluency};
\node[lbl, above=0.18cm of ling]
  {TTR, syntax, IUs,\\BERT embeddings};
\node[lbl, below=0.22cm of out]
  {\textbf{AD} / \textbf{HC}};

\end{tikzpicture}
\end{center}
\vspace{0.3em}
\end{block}

\end{column}

% ===========================================================
% COLUMN SEPARATOR
% ===========================================================
\begin{column}{\colsep}\end{column}

% ===========================================================
% RIGHT COLUMN
% ===========================================================
\begin{column}{\rightcol}

% -----------------------------------------------------------
\begin{block}{4\enspace Feature Extraction Approaches}
\justifying
Three complementary feature streams are used in the literature:

\vspace{0.5em}
\renewcommand{\arraystretch}{1.4}
\begin{tabular}{@{}>{\bfseries}p{0.22\textwidth} p{0.72\textwidth}@{}}
  \toprule
  \rowcolor{Cprimary!12}
  Stream & Key Features and Methods \\
  \midrule
  Acoustic
    & eGeMAPS (88 features: F0, energy, spectral flux, voice quality,
      MFCCs); pause duration and position; disfluency rate;
      speech rate; self-supervised representations
      (wav2vec\,2.0, Whisper) \\
  Linguistic / NLP
    & Type-token ratio (TTR), part-of-speech distribution,
      syntactic depth, information unit (IU) coverage,
      semantic coherence; fine-tuned BERT / RoBERTa
      on ASR transcripts \\
  Multimodal
    & Late fusion of acoustic + NLP (Martinc \& Pollak, 2020);
      joint BiLSTM over text, disfluency \& pause tokens
      (Rohanian et al., 2021); end-to-end WavBERT
      (Zhu et al., 2021) \\
  Graph-based
    & Spatio-semantic scene graphs capturing \emph{which} objects
      are mentioned and \emph{how} they are spatially related
      (Ambadi et al., 2021) \\
  \bottomrule
\end{tabular}
\end{block}

\vspace{\blockskip}

% -----------------------------------------------------------
\begin{block}{5\enspace Benchmark Performance}
\justifying
Results on the ADReSS / ADReSSo test sets (AD classification accuracy):

\vspace{0.5em}
\renewcommand{\arraystretch}{1.45}
\begin{tabular}{@{}p{0.44\textwidth} p{0.28\textwidth} >{\centering\arraybackslash}p{0.16\textwidth}@{}}
  \toprule
  \rowcolor{Cprimary!12}
  \textbf{System} & \textbf{Method} & \textbf{Acc.} \\
  \midrule
  Fraser et al.\ (2016)             & 370 linguistic feat.\ + SVM        & 81.9\% \\
  Luz et al.\ (2020) — baseline     & eGeMAPS acoustic                   & 62.5\% \\
  Balagopalan et al.\ (2020)        & Fine-tuned BERT                    & $\sim$83\% \\
  Martinc \& Pollak (2020)          & TF-IDF + acoustic fusion           & 81.3\% \\
  Rohanian et al.\ (2021)           & Multimodal BiLSTM                  & 84.0\% \\
  Zhu et al.\ (2021) WavBERT        & wav2vec\,2 + BERT + pause embed.   & 83.1\% \\
  \rowcolor{Cprimary!20}
  \textbf{ADReSS-M top} (2024)      & wav2vec\,2 / Whisper               & \textbf{87.0\%} \\
  \bottomrule
\end{tabular}

\vspace{0.5em}
\noindent Performance rose from 62.5\% (acoustic baseline, 2020) to
\textbf{87\%} (self-supervised end-to-end, 2024) --- a 24.5 pp gain
over four years of shared-task competition.
\end{block}

\vspace{\blockskip}

% -----------------------------------------------------------
\begin{alertblock}{6\enspace The Reality Gap}
\justifying
\begin{itemize}[leftmargin=1.4em, itemsep=0.35em, topsep=0pt]
  \item \textbf{Small datasets:} DementiaBank has $\approx$300 unique
        speakers --- far too few for robust population-level generalization
  \item \textbf{Cross-corpus drop:} ADReSS-M (Luz et al., 2024) trained
        on English, tested on Greek $\rightarrow$ baseline accuracy
        \textbf{drops to 73.9\%}
  \item \textbf{Label noise:} MMSE scores are coarse, examiner-variable
        proxies for dementia stage
  \item \textbf{Demographic under-coverage:} age, education level, accent,
        and language diversity are poorly controlled
  \item \textbf{Artifact learning:} Guo et al.\ (2021) show BERT can learn
        dataset-specific cues; Favaro et al.\ (2024) attempt to isolate
        invariant cognitive markers
\end{itemize}
\end{alertblock}

\vspace{\blockskip}

% -----------------------------------------------------------
\begin{exampleblock}{7\enspace Critical Reflection}
\justifying
\begin{itemize}[leftmargin=1.4em, itemsep=0.35em, topsep=0pt]
  \item \textbf{Accuracy $\neq$ clinical validity:} no prospective
        trials on undiagnosed community populations exist to date
  \item \textbf{Interpretability:} which speech features are
        \emph{genuinely} diagnostic vs.\ spuriously correlated with
        age, education, or recording conditions?
  \item \textbf{Ethical stakes:} false positives carry social and
        psychological consequences; fairness across demographics is
        under-studied
  \item \textbf{Path forward:} larger multilingual corpora, inherently
        interpretable architectures, and prospective clinical validation
        are the three critical bottlenecks
\end{itemize}
\end{exampleblock}

\end{column}
\end{columns}

% ===========================================================
% FULL-WIDTH CONCLUSIONS
% ===========================================================
\vspace{1.2cm}

\begin{beamercolorbox}[
  wd=\linewidth,
  sep=0.7cm,
  rounded=true,
  shadow=false
]{block body}
  {\large\bfseries\color{Cprimary} Conclusions}\quad
  Speech-based cognitive decline screening has matured from
  hand-crafted linguistic features to self-supervised end-to-end
  models, with ADReSS benchmark accuracy reaching 87\%.
  Yet cross-corpus experiments reveal substantial generalization
  failures, and no large-scale prospective clinical validation
  exists.
  Closing the gap demands diverse multilingual datasets,
  interpretable model outputs, and rigorous clinical trial design ---
  not merely higher benchmark numbers.
\end{beamercolorbox}

% ===========================================================
% REFERENCES
% ===========================================================
\vspace{0.9cm}

\begin{block}{References}
\footnotesize\setlength{\columnsep}{2cm}
\begin{multicols}{2}
\begin{itemize}[leftmargin=0pt, label={}, itemsep=0.18em]
  \item Ambadi, P.\,S., et al. (2021). Spatio-semantic graphs from
        picture description: applications to detection of cognitive
        impairment. \textit{Frontiers in Neurology}, 12, 795374.
        \url{https://doi.org/10.3389/fneur.2021.795374}

  \item Balagopalan, A., Eyre, B., Rudzicz, F., \& Novikova, J.
        (2020). To BERT or not to BERT: Comparing speech and
        language-based approaches for Alzheimer's disease detection.
        \textit{Proc.\ Interspeech 2020}, 2167--2171.
        \url{https://doi.org/10.21437/Interspeech.2020-2557}

  \item Becker, J.\,T., et al. (1994). The natural history of
        Alzheimer's disease. \textit{Archives of Neurology},
        51(6), 585--594.
        \url{https://doi.org/10.1001/archneur.1994.00540180063015}

  \item Favaro, A., et al. (2024). Discovering invariant patterns of
        cognitive decline via an automated analysis of the Cookie
        Thief picture description task. \textit{Proc.\ Odyssey 2024},
        201--208.
        \url{https://doi.org/10.21437/Odyssey.2024-29}

  \item Fraser, K.\,C., Meltzer, J.\,A., \& Rudzicz, F. (2016).
        Linguistic features identify Alzheimer's disease in narrative
        speech. \textit{Journal of Alzheimer's Disease}, 49(2),
        407--422.
        \url{https://doi.org/10.3233/JAD-150520}

  \item Guo, Y., et al. (2021). Crossing the ``Cookie Theft'' corpus
        chasm: Applying what BERT learns from outside data to the
        ADReSS challenge dementia detection task.
        \textit{Frontiers in Computer Science}, 3, 642517.
        \url{https://doi.org/10.3389/fcomp.2021.642517}

  \item Luz, S., et al. (2020). Alzheimer's dementia recognition
        through spontaneous speech: The ADReSS challenge.
        \textit{Proc.\ Interspeech 2020}, 2172--2176.
        \url{https://doi.org/10.21437/Interspeech.2020-2571}

  \item Luz, S., et al. (2021). Detecting cognitive decline using
        speech only: The ADReSSo challenge.
        \textit{Proc.\ Interspeech 2021}, 3780--3784.
        \url{https://doi.org/10.21437/Interspeech.2021-1220}

  \item Luz, S., et al. (2024). An overview of the ADReSS-M Signal
        Processing Grand Challenge on Multilingual Alzheimer's
        Dementia Recognition Through Spontaneous Speech.
        \textit{IEEE Open Journal of Signal Processing}, 5, 738--749.
        \url{https://doi.org/10.1109/ojsp.2024.3378595}

  \item Rohanian, M., Hough, J., \& Purver, M. (2021). Alzheimer's
        dementia recognition using acoustic, lexical, disfluency and
        speech pause features robust to noisy inputs.
        \textit{Proc.\ Interspeech 2021}, 3820--3824.
        \url{https://doi.org/10.21437/Interspeech.2021-1633}

  \item Zhu, Y., et al. (2021). WavBERT: Exploiting semantic and
        non-semantic speech using Wav2vec and BERT for dementia
        detection. \textit{Proc.\ Interspeech 2021}.
        \url{https://doi.org/10.21437/Interspeech.2021-332}
\end{itemize}
\end{multicols}
\end{block}

\end{frame}
\end{document}
